\chapter*{专题练习4}
填空题：
\begin{enumerate}
    \item 已知函数$f(x)$是定义在$(0, +\infty)$上的单调函数，且对$x \in (0, +\infty)$都有$f(f(x) - \dfrac{4}{x}) = 4$，则$f(x) = $\blkda{$2+\frac{4}{x}$}
    \begin{jieda}
        由于函数$f(x)$是定义在$(0, +\infty)$的单调函数，则必有$f(x) - \dfrac{4}{x}$为常数才可以使得$f(f(x) - \dfrac{4}{x}) = 4$，故可以设$f(x) = c + \dfrac{4}{x}$，因此有$f(c) = c + \dfrac{4}{c} = 4$，解得$c = 2$，因此$f(x) = 2+\dfrac{4}{x}$
    \end{jieda}
    \item 已知函数$f(x) = x^3+x$，若$f(t^2+2t+k)+f(-2t^2+2t-5)>0$对任意的$t \in [0, 5]$都成立，则$k$的取值范围是\blkda{$(10, +\infty)$}.
    \begin{jieda}
        $f(x)$为奇函数，在$\mathbf{R}$上单调增，根据对称性可以将函数值上的不等关系转化为变量的不等关系，即$(t^2+2t+k)+(-2t^2+2t-5)>0$在$t \in [0, 5]$恒成立.\\
        参变分离有$t^2-4t+5 < k$在$t \in [0, 5]$恒成立，因此有$k > 10$ \\
        \ps $f(x)$关于$(a, b)$对称，则$f(x_1) + f(x_2) = 2b \Leftrightarrow x_1 + x_2 = 2a$ \\
        $f(x)$关于$(a, b)$对称，且在定义域内单调增，则$f(x_1) + f(x_2) < 2b \Leftrightarrow x_1 + x_2 < 2a$ \\
        $f(x)$关于$(a, b)$对称，且在定义域内单调增，则$f(x_1) + f(x_2) > 2b \Leftrightarrow x_1 + x_2 > 2a$ \\
        $f(x)$关于$(a, b)$对称，且在定义域内单调减，则$f(x_1) + f(x_2) > 2b \Leftrightarrow x_1 + x_2 < 2a$ \\
        $f(x)$关于$(a, b)$对称，且在定义域内单调减，则$f(x_1) + f(x_2) < 2b \Leftrightarrow x_1 + x_2 > 2a$ \\
        以上结论可以通过代数推导，更推荐结合图像理解.
    \end{jieda}
    \item 已知函数$f(x) = 2^x - (\frac{1}{2})^x+3$，且$f(m^2)+f(-5m+4)<6$，则实数$m$的取值范围是\blkda{$(1, 4)$}.
    \begin{jieda}
        直接应用上面结论即可，也可以设$g(x) = f(x) - 3$，将$f(m^2)+f(-5m+4)<6$转化为$g(m^2)+g(-5m+4)<0$后按照更熟悉的奇函数形式处理.
    \end{jieda}
    \item 已知函数$y = f(3x+1)$为偶函数，且在$[0, +\infty)$上为严格增函数，若$f(x) < f(2x+1)$，则$x$的取值范围是\blkda[5]{$(-\infty, -1) \cup (\frac{1}{3}, +\infty)$}.
    \begin{jieda}
        函数$y = f(3x+1)$为偶函数，且在$[0, +\infty)$上为严格增函数即$f(x)$关于$x = 1$对称，且“开口向上”，因此离对称轴$x = 1$近的变量的函数值小，由此$f(x) < f(2x+1)$可以转化为$|x-1| < |2x|$，两侧平方有$(x-1)^2 < 4x^2$，解得$x \in (-\infty, -1) \cup (\frac{1}{3}, +\infty)$
    \end{jieda}
    \item 已知函数$f(x)$在定义域$[2-a, 3]$上是偶函数，在$[0, 3]$上严格单调递减，并且$f(-m^2-\frac{a}{5}) > f(-m^2+2m-2)$，则$m$的取值范围是\blkda[3]{$[1-\sqrt{2}, \frac{1}{2})$}.
    \begin{jieda}
        首先根据定义域的对称性可以求得$a = 5$，进一步根据对称性可以判断出函数“开口向下”，离$y$轴越近函数值越大，再结合定义域的要求，则$f(-m^2-\frac{a}{5}) > f(-m^2+2m-2) \Leftrightarrow \left \{ \begin{array}{l}
             |-m^2-1| < |-m^2+2m-2|  \\
             -3 \leq -m^2-1 \leq 3  \\
             -3 \leq -m^2+2m-2 \leq 3 
        \end{array}\right.$第一个配方后可以直接去绝对值，最终得到$m \in [1-\sqrt{2}, \frac{1}{2})$ \\
        \ps \textbf{奇函数和偶函数的定义域都是关于原点对称的，若要证明奇偶性首先要判断定义域！}
    \end{jieda}
    \item 已知函数$f(x) = \left\{  
     \begin{array}{ll}  
        x^2+x, & 0 < x < 2, \\  
        -2x+8, & x \geq 2.
     \end{array}
     \right.$ 若$f(a) = f(a+2)$, 则$f(\frac{1}{a})$的值是\blkda{2}
     \begin{jieda}
         根据题意可知$0 < a < 2$，否则$a > 2, a+2 > 2$，由单调性可知无法满足$f(a) = f(a+2)$,因此有$a^2 + a = -2(a+2)+8$。解得$a = 1$(舍去-4)，故$f(\dfrac{1}{a}) = f(1) = 2$
     \end{jieda}
    \item 设函数$f(x) = x - \dfrac{1}{x}$。对任意$x \in [1, +\infty), f(mx) + mf(x) < 0$恒成立，则实数$m$的取值范围是\blkda{$(-\infty, -1)$}.
    \begin{jieda}
        $f(x)$为奇函数，在$(-\infty, 0), (0, +\infty)$内单调递增，$f(1) = f(-1) = 0$。
        \begin{enumerate}
            \item 若$m > 0$，则$x \in [1, +\infty)$时$f(mx) > 0, mf(x) > 0$，不满足题意。
            \item 若$m = 0$，$mx$不在$f(x)$定义域内。
            \item 若$m < 0$，则$f(mx)+mf(x) = mx - \dfrac{1}{mx} + mx - \dfrac{m}{x} = 2mx - \dfrac{1+m^2}{mx}  < 0$，整理得$2mx < \dfrac{1+m^2}{mx} \Rightarrow 2m^2x^2 > 1+m^2 \Rightarrow m^2 > \dfrac{1}{(2x^2-1)}$。因为当$x \in [1, +\infty)$时，$\dfrac{1}{(2x^2-1)} \leq 1$，故$m < -1$
        \end{enumerate}
    \end{jieda}
    \item 函数$y=\frac{3-x}{1+2x}$的严格递减区间是\blkda[4]{$(-\infty,-\frac{1}{2})$和$(-\frac{1}{2},+\infty)$}, 对称中心是\blkda{$(-\frac{1}{2},-\frac{1}{2})$}, \\
    值域是\blkda[2]{$\{t\mid t\neq-\frac{1}{2}\}$}. 
    \begin{jieda}
        反比例函数对称中心的横坐标是分母为零的值，纵坐标是分离出的常数.
    \end{jieda}
    \item 已知函数$f(x)=x^5+ax^3+bx-4$, 且$f(-2)=10$, 则$f(2)$=\blkda{$-18$}
    \begin{jieda}
    $f(x)+f(-x)=-8$，函数关于$(0, -4)$中心对称。\\
    \ps 设$g(x) = f(x)+4$，则$g(x)$为奇函数. (多项式函数是奇函数等价于偶次方的系数为零, 是偶函数等价于奇次方的系数为零.)
    \end{jieda}
    \item 已知$f(x) = \frac{(2^x+1)^2}{2^x \cdot x}+1$在$[-2018, -1]\cup[1, 2018]$上最大值为$M$,最小值为$m$,则$M+m = $\blkda{$2$}.
    \begin{jieda}
    设$g(x) = \frac{(2^x+1)^2}{2^x \cdot x}$，则$g(x) = \frac{2^x+\frac{1}{2^x}+2}{x} = \frac{2^x+2^{-x}+2}{x}$，故$g(-x) = -g(x)$，即$g(x)$为奇函数，故$M+m = 2$ \\
    \ps \textbf{涉及最大值与最小值和的题目一般可以优先考虑对称性}.
    \end{jieda}
    \item 已知函数$f(x) = \frac{3^{x+1}+1}{3^x+1}$在$[-2021, 2021]$上的最大值和最小值分别为$M, m$，则$M+m = $\blkda{$4$}.
    \begin{jieda}
        $g(x) = f(x) - 2 = \frac{3^x - 1}{3^x+1}$为奇函数，故$f(x)$关于$(0, 2)$中心对称，$M+m = 4$ \\
        \ps 本题如果没有发现$f(x) - 2$是奇函数，利用单调性也是可以解决的，计算量相对大. \\
        \ps 其实$y=\frac{a^x-1}{a^x+1}$, $y=\frac{a^x+1}{a^x-1}$都是奇函数.
    \end{jieda}
    \item 设函数$f(x)$的定义域为$\mathbf{R}$，$f(x+1)$是奇函数，$f(x+2)$是偶函数，当$x \in [1, 2]$时，$f(x) = kx+b$，若$f(0) + f(3) = 8$，则$f(\frac{2023}{2}) = $\blkda{4}
    \begin{jieda}
    因为$f(x)$关于$(1, 0)$中心对称，当$x \in [1, 2]$时，$f(x) = kx+b$为连续函数，所以$f(1) = 0$ 且在$[0, 2]$上$f(x) = kx+b$. 又因为$f(x)$关于$x = 2$对称，所以$f(3) = f(1)$, 故$f(0) = 8, f(\frac{1}{2}) = 4, f(\frac{3}{2}) = -4$，即$f(\frac{8k+1}{2}) = 4, f(\frac{8k+3}{2}) = -4, f(\frac{8k+5}{2}) = -4, f(\frac{8k+7}{2}) = 4 (k \in \mathbf{Z})$因此$f(\frac{2023}{2}) = 4$ \\
    \ps 如果对称中心的横坐标在定义域内，则对称中心在函数图象上，如果对称中心横坐标不在定义域内，则对称中心不在函数图象上；
    \end{jieda}
    \item 已知函数$g(x) = x^3 - 9x^2 + 29x - 30$，$g(m) = -12, g(n) = 18$，则$m+n = $\blkda{6}
    \begin{jieda}
        对于三次函数$y = ax^3+bx^2+cx+d$，当$b^2 - 3ac \leq 0$时，其为单调函数，因此$g(x)$是单调函数. \\
        对于三次函数$y = ax^3+bx^2+cx+d$，其对称中心为$(-\frac{b}{3a}, f(-\frac{b}{3a}))$，故$g(x)$对称中心为$(3, 3)$，因为$\frac{g(m) + g(n)}{2} = 3$，所以$m+n = 6$
    \end{jieda}
    \item 已知函数$f(x) = \left \{ \begin{array}{ll}
        1 - |x| & x \leq 1 \\
        (x-1)^2 & x > 1
    \end{array}\right.$，函数$g(x) = b - f(1-x)$，如果$y = f(x) - g(x)$恰有两个零点, 则实数$b$的取值范围是\blkda[3]{$\{\frac{3}{4}\} \cup (1, +\infty)$}
    \begin{jieda}
        设$h(x) = f(x) - g(x) = f(x) + f(1-x) - b$，$h(x)$的零点即$y = f(x) + f(1-x)$与$y = b$的交点的横坐标. 因为$f(x)$与$f(1-x)$关于$x = \frac{1}{2}$对称，所以$y = f(x) + f(1-x)$关于$x = \frac{1}{2}$对称. 当$x \in [0, 1]$时，$f(x) + f(1-x) = 1$；当$x > 1$或$x < 0$时，$f(x) + f(1-x)$为二次函数，最小值为$\dfrac{3}{4}$ \\
        根据对称性绘制图像可知当$b \in \{\frac{3}{4}\} \cup (1, +\infty)$时满足要求. \\
        \ps \textbf{两个互对称的函数相加的结果会具有自对称性.}
    \end{jieda}
    \item 已知定义在$\mathbf{R}$上的偶函数$f(x)$满足$f(x) = -f(\pi - x)$，当$x \in [0, \frac{\pi}{2})$时$f(x) = \sqrt{x}$，则函数$y = f(x)$在$[-\frac{3\pi}{2}, \frac{5\pi}{2}]$上的零点个数为\blkda{9}
        \begin{jieda}
            $f(x) = -f(\pi - x)$即$f(x)$关于$(\frac{\pi}{2}, 0)$中心对称，又因为$f(x)$是偶函数，所以$2\pi$是$f(x)$的一个周期. 因为定义域为$\mathbf{R}$，所以$f(\frac{k\pi}{2}) = 0, k \in \mathbf{Z}$；结合图像可知有9个零点 \\
        \ps 两个横坐标不同位置的对称性会产生周期性
        \begin{enumerate}
            \item 两个对称轴的横坐标之差的二倍是函数的周期：$f(x) = f(2a-x), f(x) = f(2b - x) \Leftrightarrow 2|a-b|$是$f(x)$的一个周期。
            \item 两个纵坐标相同的对称中心的横坐标之差的二倍是函数的周期：$f(x) + f(2a-x) = 2m, f(x) + f(2b - x) = 2m \Leftrightarrow 2|a-b|$是$f(x)$的一个周期。
            \item 一个对称轴和一个对称中心的横坐标之差的四倍是函数的周期：$f(x) = f(2a-x), f(x) + f(2b - x) = 2m \Leftrightarrow 4|a-b|$是$f(x)$的一个周期。
        \end{enumerate}
        \end{jieda}
        \item 已知函数$f(x)$的定义域为$\mathbf{R}$，$f(x+y)+f(x-y) = f(x) \cdot f(y)$，$f(1) = 0$，则$f(2022) = $\blkda{$-2 \mbox{或} 0$}
        \begin{jieda}
            令$y = 1$，则对于任意的$x \in \mathbf{R}$有$f(x+1) + f(x-1) = 0$ \\ 故迭代有$f(x+3) + f(x+1) = 0$，因此有$f(x + 3) = f(x-1)$，即$4$是$f(x)$的一个周期. 故$f(2022) = f(2)$\\
            代入$x = 0, y = 0$，有$f(0) + f(0) = f(0) \cdot f(0)$，故$f(0) = 0$或$f(0) = 2$ \\
            代入$x = 1, y = 1$，由$f(2) + f(0) = f(1) \cdot f(1)$，故$f(2) = -f(0)$ \\
            因此$f(2022) = -2 \mbox{或} 0$
        \end{jieda}
        \item $f(x)$是定义在$\mathbf{R}$上的以3为周期的奇函数，且$f(2) = 0$，则方程$f(x) = 0$在区间$[0, 6]$内的解的个数的最小值是\blkda{$9$}. 
        \begin{jieda}
            $f(x)$为$\mathbf{R}$上奇函数，故$f(0) = 0$，又因为$f(2) = 0$，由周期性$f(-1) = 0$，由奇函数性质进一步可以推得$f(-2) = 0$，故由周期性可知$f(n) = 0 (n \in \mathbf{Z})$ \\
            由奇函数和周期性可知$f(1.5) = f(-1.5) = 0$，进一步可知$f(4.5) = 0$ \\
            \ps 对于周期为$T$的奇函数，必有$f(\frac{T}{2}) = f(-\frac{T}{2}) = 0$
        \end{jieda}
        \item 已知定义在$\mathbf{R}$上的奇函数$f(x)$满足$f(x+1) = f(1-x)$，且当$x \in [0, 1]$时，$f(x) = 2^x - m$，则$f(2023) =$ \blkda{-1}
        \begin{jieda}
            根据$f(x+1) = f(1-x)$可知$f(x)$关于$x = 1$轴对称，又因为$f(x)$为奇函数，因此$f(x)$的周期为4，故$f(2023) = f(-1) = -f(1)$。\\
            因为$f(x)$为奇函数，所以$f(0) = 0$，因此$m = 1$，所以$f(1) = 2 - 1 = 1$故$f(2023) = -1$
        \end{jieda}
    \end{enumerate}
\clearpage
选择题：
\begin{enumerate}
    \item 已知函数$f(x)$和$g(x)$均为奇函数，$h(x) = a \cdot f(x) - b \cdot g(-x) + 8$在区间$(0, +\infty)$上有最大值9，那么$h(x)$在$(-\infty, 0)$上的最小值为\dotfill(\kh{C})
    \xx{-9}{-1}{7}{1}
    \begin{jieda}
        由于$g(x)$为奇函数，故$h(x) = af(x) - bg(x) + 8$, 设$h_1(x) = h(x) - 8 = af(x) - bg(x)$，易知$h_1(x)$在区间$(0, +\infty)$上有最大值1，且$h_1(x)$为奇函数，所以$h_1(x)$在区间$(-\infty， 0)$上有最小值-1。因此$h(x)$在区间$(-\infty， 0)$上最小值为7
    \end{jieda}
    \item 定义在$\bR$上的偶函数$f(x)$满足: 对任意的$x_1,x_2\in[0,+\infty)$, $x_1\neq x_2$, 有$\frac{f(x_2)-f(x_1)}{x_2-x_1}<0$, 则\dotfill(\kh{A})
        \xx{$f(3)<f(-2)<f(1)$}{$f(1)<f(-2)<f(3)$}{$f(-2)<f(1)<f(3)$}{$f(3)<f(1)<f(-2)$}
        \begin{jieda}
        等价定义: 
        \begin{itemize}
        \item 严格单调增: $\frac{f(x_2)-f(x_1)}{x_2-x_1}>0(x_1\neq x_2)$; 
        \item 严格单调减: $\frac{f(x_2)-f(x_1)}{x_2-x_1}<0(x_1\neq x_2)$. 
        \end{itemize}
    \end{jieda}
    \item 已知函数$f(x)\neq -1$, 且对定义域内任意的$x$总有关系$(f(x+\pi)+1)(f(x)+1)=2$, 那么下列结论中正确的是\dotfill(\kh{B})
    \xx{$f(x)$是周期为$\pi$的周期函数}
    {$f(x)$是周期为$2\pi$的周期函数}
    {$f(x)$是周期为$\frac{\pi}{2}$的周期函数}
    {$f(x)$不是周期函数}
    \begin{jieda}
    由条件可得$(f(x+2\pi)+1)(f(x+\pi)+1)=2$, 于是$f(x+2\pi)=f(x)$. \\
    \ps $f(x) + f(x+t) = c\Leftrightarrow 2t$是$f(x)$的一个周期。\\
    进阶版：$f(x+a) + f(x+b) = c\Leftrightarrow 2|a-b|$是$f(x)$的一个周期。 \\
    $f(x) \cdot f(x+t) = c (c \neq 0)\Leftrightarrow 2t$是$f(x)$的一个周期。\\
    进阶版：$f(x+a) \cdot f(x+b) = c (c \neq 0)\Leftrightarrow 2|a-b|$是$f(x)$的一个周期。
    \end{jieda}
    \item 已知$f(x)$是定义域为$(-\infty, +\infty)$的奇函数，满足$f(1-x) = f(1+x)$。若$f(1) = 2$，则$f(1)+f(2)+f(3)+...+f(50) = $\dotfill(\kh{C})
    \xx{-50}
    {0}
    {2}
    {50}
    \begin{jieda}
        因为$f(x)$为奇函数，且$f(1-x) = f(1+x)$，故$f(x)$具有周期性，最小正周期为4，即$f(x) = f(x+4)$. 由题意可知$f(0) = 0, f(1) = 2$, 可以推出$f(2) = 0, f(3) = f(-1) = -2$，由此可知$f(4k)+f(4k+1)+f(4k+2)+f(4k+3)=0, k \in \mathbf{Z}$。故原式$= f(1) + f(2) = 2$ \\
        \ps 具有周期性的求和问题，先求一个周期内的和.
    \end{jieda}
\end{enumerate}
解答题：
\begin{enumerate}
    \item 函数$f(x) = \frac{x}{ax+b}$（$a, b$是非零常数），满足$f(2) = 1$，且方程$f(x) = x$仅有一个解.
    \begin{enumerate}
        \item 求$a, b$的值
        \item 是否存在常数$m$，使得对于定义域内任意的$x$，$f(x) + f(m-x) = 4$恒成立
        \item 在平面直角坐标系中，求定点$A(-3, 1)$到此函数上任意一点$P$的距离$|AP|$的最小值.
    \end{enumerate}
    \begin{jieda}
        \begin{enumerate*}
            \item 由$f(2) = 1$，$\frac{2}{2a+b} = 1$，故$2a+b = 2$；$f(x) = x$即$\frac{x}{ax+b} = x$，即$(\frac{1}{ax+b}-1)x = 0$，容易判断其有两个解$x_1  = 0, x_2 = \frac{1-b}{a}$.要符合题意只可能是$b = 1$，两根重合.进一步推得$a = \frac{1}{2}$ \\
            \item $f(x) = \frac{2x}{x+2}$，为反比例型函数，其对称中心为$(-2, 2)$，因此$f(x)$满足$f(x) + f(-4-x) = 4$. 由此可知$m = -4$ \\
            \item 设$P(x_0, \frac{2x_0}{x_0+2})$，故$|AP| = \sqrt{(x_0+3)^2+(\frac{2x_0}{x_0+2}-1)^2}   = \sqrt{(x_0+3)^2+(1-\frac{4}{x_0+2})^2}$. 设$t = x_0+2 (t \neq 0)$，则$|AP| = \sqrt{(t+1)^2+(1-\frac{4}{t})^2} = \sqrt{(t-\frac{4}{t})^2+2(t-\frac{4}{t})+10}$，设$k = t - \frac{4}{t}, k \in \mathbf{R}$，$|AP| = \sqrt{k^2+2k+10} = \sqrt{(k+1)^2+9} \geq 3$（当$x = \frac{-5\pm\sqrt{17}}{2}$时取得最小值）
        \end{enumerate*}
    \end{jieda}
    \begin{vspace_}
    \vsp[5]
    \end{vspace_}
    \item 已知$f(x)$是$\bR$上的函数, 且$f(x+2)(1-f(x))=1+f(x)$. 
        \begin{enumerate*}[1.]
        \item 求证: $f(x)$是周期函数;
        \item 若$f(4)=-\sqrt{3}$, 求$f(2020)$. 
        \end{enumerate*} 
        \begin{jieda}
        \begin{enumerate}[1.]
        \item 若$f(x)$为常数函数, 则设为$f(x)=t$, 于是$t(1-t)=1+t$, 无解, 故$f(x)$不是常函数. 
        \begin{dingautolist}{172}
            \item 若$f(x)=1$有解, 设为$x_0$, 则$f(x_0+2)(1-f(x_0))=1+f(x_0)$, 故$f(x_0)=-1$. 矛盾. \\
            若$f(x)=0$有解, 设为$x_0$, 则$f(x_0+2)(1-f(x_0))=1+f(x_0)$, 故$f(x_0+2)=1$. 显然不可能(见上一行).
            \item 由(i)知$f(x)=1$无解, 且$f(x)=0$无解, \\
            则$f(x+4)=\frac{1+f(x+2)}{1-f(x+2)}=\frac{1+\frac{1+f(x)}{1-f(x)}}{1-\frac{1+f(x)}{1-f(x)}}=-\frac{1}{f(x)}$, 故$f(x+8) = -\frac{1}{f(x+4)} = f(x)$.
        \end{dingautolist}
        综上, 8是$f(x)$的一个周期.
        \item $f(2020)=f(4)=-\sqrt{3}$. 
        \end{enumerate}
        \end{jieda}
    \begin{vspace_}
    \vsp[5]
    \end{vspace_}
    \item 已知函数$f(x)$是周期为$4$的函数，若$x \in [-2, 2]$时，$f(x) = x^2 + x^4$，则$x \in [2, 6]$时求$f(x)$解析式.
    \begin{jieda}
        任取$x \in [2, 6]$，因为$f(x)$周期为4，所以$f(x) = f(x - 4)$，又因为$x - 4 \in [-2, 2]$，故$f(x-4) = (x-4)^2 + (x-4)^4$，因此$f(x) = (x-4)^2 + (x-4)^4 (x \in [2, 6])$ 
    \end{jieda}
    \begin{vspace_}
    \vsp[5]
    \end{vspace_}
    \item 已知$y=f(x)(x\zr)$是奇函数, 若当$x<0$时, $f(x)=x^{\frac{1}{3}}+2^x-1$, 求函数$f(x)$的解析式.
    \begin{jieda}
    当$x>0$时, $f(x)=-f(-x)=-((-x)^{\frac{1}{3}}+2^{-x}-1)=x^{\frac{1}{3}}-2^{-x}+1$. \\
    又上述两式当$x=0$时也满足$f(0)=0$.  \\
    故$f(x) = \begin{cases}
    x^{\frac{1}{3}}+2^x-1, x\leq 0\\
    x^{\frac{1}{3}}-2^{-x}+1, x>0
    \end{cases}$. 
    \end{jieda}
    \begin{vspace_}
    \vsp[5]
    \end{vspace_}
    \item 已知函数$f(x)$在$\mathbf{R}$上有定义，对任意实数$a>0$和任意实数$x$，都有$f(ax) = af(x)$.
    \begin{enumerate}
    \item 证明$f(0) = 0$
    \item 证明$f(x) = \left\{  
         \begin{array}{ll}  
            kx, & x \geq 0, \\  
            hx & x < 0.
         \end{array}
    \right.$ 其中$k, h$均为常数。
    \item 当上一问中的$k > 0$，设$g(x) = \dfrac{1}{f(x)} + f(x) (x > 0)$，讨论$g(x)$在$(0, +\infty)$的单调性并求最值. 
    \end{enumerate}    
    \begin{jieda}
        \begin{enumerate}
        \item 令$x = 0$，有$f(0) = a f(0), a > 0$，故$f(0) = 0$
        \item 当$x > 0$时，令$x = 1$，$f(a) = af(1)$，因为对任意$a$均成立，故可将$a$作为自变量，写为$x$，有$f(x) = f(1) \cdot x$，取$k = f(1)$，有$f(x) = kx, x \geq 0$；当$x < 0$时，令$x = -1$，$f(-a) = -af(-1)$，故同理可以设$h = -f(-1)$，有$f(x) = hx, x < 0$。
        \item $g(x) = kx + \dfrac{1}{kx} x > 0, k > 0$，换元$u = kx$，设$g_1(u) = u+\dfrac{1}{u}$，根据对勾函数性质可知$g_1(u)$在$(0, 1)$单调递减，在$[1, +\infty)$单调递增，因此有$g(x)$在$(0, \dfrac{1}{k})$单调递减，在$[\dfrac{1}{k}, +\infty)$单调递增。其最小值为2。
\end{enumerate}
    \end{jieda}
\end{enumerate}